% Unidad R572: EGA I, capítulo 0, (1.0.4)--(1.0.7).
% Traducción: OpenAI Codex, GPT-6.1 Sol, esfuerzo Ultra.
\begin{env}[1.0.4]
\label{0.1.0.4-fr}
Si $A$ es un anillo (conmutativo) y $B$ un anillo no necesariamente conmutativo, dar una estructura de \emph{$A$-álgebra} sobre $B$ equivale a dar un homomorfismo de anillos $\varphi:A\to B$ tal que $\varphi(A)$ esté contenido en el centro de $B$.
Para todo ideal $\mathfrak{J}$ de $A$, $\mathfrak{J}B=B\mathfrak{J}$ es entonces un ideal bilateral de $B$, y para todo $B$-módulo $M$, $\mathfrak{J}M$ es entonces un $B$-módulo igual a $(B\mathfrak{J})M$.
\end{env}

\begin{env}[1.0.5]
\label{0.1.0.5-fr}
No volveremos sobre las nociones de \emph{módulo de tipo finito} y de \emph{álgebra} (conmutativa) \emph{de tipo finito}; decir que un $A$-módulo $M$ es de tipo finito significa que existe

\oldpage[0\textsubscript{I}]{12}
una sucesión exacta $A^p\to M\to 0$.
Se dice que un $A$-módulo $M$ admite una \emph{presentación finita} si es isomorfo al conúcleo de un homomorfismo $A^p\to A^q$; en otras palabras, si existe una sucesión exacta $A^p\to A^q\to M\to 0$.
Se observará que, sobre un anillo \emph{noetheriano} $A$, todo $A$-módulo de tipo finito admite una presentación finita.

Recordemos que una $A$-álgebra $B$ se dice \emph{entera sobre $A$} si todo elemento de $B$ es raíz en $B$ de un polinomio mónico con coeficientes en $A$; equivale a decir que todo elemento de $B$ está contenido en una subálgebra de $B$ que es un \emph{$A$-módulo de tipo finito}.
Cuando esto ocurre y $B$ es conmutativa, la subálgebra de $B$ generada por una parte finita de $B$ es un $A$-módulo de tipo finito; por tanto, para que el álgebra conmutativa $B$ sea entera y de tipo finito sobre $A$, es necesario y suficiente que $B$ sea un $A$-módulo de tipo finito; se dice entonces también que $B$ es una $A$-álgebra \emph{entera finita} (o simplemente \emph{finita}, si no hay riesgo de confusión).
Se observará que, en estas definiciones, no se supone que el homomorfismo $A\to B$ que define la estructura de $A$-álgebra sea inyectivo.
\end{env}

\begin{env}[1.0.6]
\label{0.1.0.6-fr}
Un anillo \emph{íntegro} es un anillo en el que el producto de una familia finita de elementos $\neq 0$ es $\neq 0$; equivale a decir que, en un anillo de este tipo, se tiene $0\neq 1$ y que el producto de dos elementos $\neq 0$ no es nulo.
Un ideal \emph{primo} de un anillo $A$ es un ideal $\mathfrak{p}$ tal que $A/\mathfrak{p}$ sea íntegro; esto implica, por tanto, que $\mathfrak{p}\neq A$.
Para que un anillo $A$ tenga al menos un ideal primo, es necesario y suficiente que $A\neq\{0\}$.
\end{env}

\begin{env}[1.0.7]
\label{0.1.0.7-fr}
Un anillo \emph{local} es un anillo $A$ en el que existe un único ideal maximal; este ideal es entonces el complementario de los elementos invertibles y contiene todos los ideales $\neq A$.
Si $A$ y $B$ son dos anillos locales, y $\mathfrak{m}$ y $\mathfrak{n}$ sus respectivos ideales maximales, se dice que un homomorfismo $\varphi:A\to B$ es \emph{local} si $\varphi(\mathfrak{m})\subset\mathfrak{n}$ (o, equivalentemente, si $\varphi^{-1}(\mathfrak{n})=\mathfrak{m}$).
Al pasar a los cocientes, un homomorfismo de este tipo define un monomorfismo del cuerpo residual $A/\mathfrak{m}$ al cuerpo residual $B/\mathfrak{n}$.
La composición de dos homomorfismos locales es un homomorfismo local.
\end{env}
